% NexusMORE independently authored Chinese research baseline 1.0.0.
% Engineering example only; not official rules, an approved paper or a result.
% Local fonts/simsun.ttc is required and is NOT distributed with this template.
\documentclass[UTF8,fontset=none,12pt,a4paper]{ctexart}
\usepackage{fontspec,geometry,amsmath,amssymb,booktabs,graphicx,longtable}
\usepackage[hidelinks]{hyperref}
\geometry{left=2.5cm,right=2.5cm,top=2.5cm,bottom=2.5cm}
\setmainfont[Path=fonts/,AutoFakeBold=2,AutoFakeSlant=0.2,SlantedFont=simsun.ttc,SlantedFeatures={FakeSlant=0.2}]{simsun.ttc}
\setCJKmainfont[Path=fonts/,AutoFakeBold=2,AutoFakeSlant=0.2]{simsun.ttc}
\setCJKsansfont[Path=fonts/,AutoFakeBold=2]{simsun.ttc}
\setCJKmonofont[Path=fonts/]{simsun.ttc}
\newcommand{\keywords}[1]{\par\noindent\textbf{关键词：}#1}
\title{LATEX_TEMPLATE_TITLE}
\author{}
\date{}
\linespread{1.3}
\setlength{\parindent}{2em}
\raggedbottom
\begin{document}
\hypersetup{pageanchor=false}
\maketitle
\thispagestyle{empty}
\begin{abstract}
本文件用于检查中文建模论文模板的章节、字体、数学表达与分页。正文中的说明段落和代数示例只用于工程排版核验，没有使用真实赛题数据，不构成求解结果或论文批准。正式论文必须从已确认的模型与结果出发，先提交中文审阅稿并分别确认正文与最终配图，再进入受控同源导出。标题、摘要和关键词应在首页完整显示，正文从第二页开始。
\keywords{中文建模；研究模板；排版核验}
\end{abstract}
\clearpage
\pagenumbering{arabic}
\hypersetup{pageanchor=true}
\section{问题重述}
实际写作应概括目标、输入和任务边界。本工程文档仅用于验证版式，不能替代真实问题描述。
\section{问题分析}
解释问题结构、候选方法及取舍依据，不把算法名称当作模型，也不将示例段落带入正式论文。
\section{模型假设}
逐项说明必要假设的依据和影响；尚未核对的条件应返回核实，不能靠假设保证结论。
\section{符号说明}
全文只使用表\ref{tab:symbols}这一张逻辑符号表；题目范围在同一表中区分，不再分设小题符号表。
\begin{longtable}{llll}
\caption{工程代数示例的统一符号}\label{tab:symbols}\\
\toprule
符号 & 含义 & 单位 & 适用范围 \\
\midrule
\endfirsthead
\toprule
符号 & 含义 & 单位 & 适用范围 \\
\midrule
\endhead
$x$ & 无量纲自变量 & $1$ & 代数排版示例 \\
$f(x)$ & 示例函数值 & $1$ & 代数排版示例 \\
\bottomrule
\end{longtable}
\section{模型建立}
下式仅用于检查原生公式和交叉引用，不代表已确认模型：
\begin{equation}\label{eq:example}
f(x)=x^2+1,\qquad x\in\mathbb{R}.
\end{equation}
\section{模型求解}
式\eqref{eq:example}可以直接代数代入；真实任务应说明Python实现、参数、停止条件和求解状态，不伪造迭代过程。
\section{结果分析}
这里不提供赛题结果。正式稿应以已确认数值和图表逐项回答问题，单位、有效位数和比较范围应清楚。
\section{误差分析与模型检验}
按题型核对约束、残差、收敛与敏感性，区分模型误差和数值误差。排版示例不替代真实检验。
\section{模型评价与推广}
优点和局限须有证据支持；不将模板成功编译解释为模型正确、官方合规或全部导出接通。
\section{结论}
正式写作应汇总已确认的回答，不引入未经验证的新结果。当前文件仅用于工程校验。
\section{参考文献}
本工程文档没有外部学术引用，不编造条目。正式稿使用实际读取并核验的编号参考文献，正文逐项对应。
\end{document}
